% Keep this generated LaTeX guide synchronized with LoadBalancer.md.
% Build from this directory: pdflatex LoadBalancer.tex (run twice).
% Options for packages loaded elsewhere
\PassOptionsToPackage{unicode}{hyperref}
\PassOptionsToPackage{hyphens}{url}
\PassOptionsToPackage{dvipsnames,svgnames,x11names}{xcolor}
%
\documentclass[
  11pt,
  a4paper,
]{article}
\usepackage{amsmath,amssymb}
\usepackage{setspace}
\usepackage{iftex}
\ifPDFTeX
  \usepackage[T1]{fontenc}
  \usepackage[utf8]{inputenc}
  \usepackage{textcomp} % provide euro and other symbols
\else % if luatex or xetex
  \usepackage{unicode-math} % this also loads fontspec
  \defaultfontfeatures{Scale=MatchLowercase}
  \defaultfontfeatures[\rmfamily]{Ligatures=TeX,Scale=1}
\fi
\usepackage{lmodern}
\ifPDFTeX\else
  % xetex/luatex font selection
\fi
% Use upquote if available, for straight quotes in verbatim environments
\IfFileExists{upquote.sty}{\usepackage{upquote}}{}
\IfFileExists{microtype.sty}{% use microtype if available
  \usepackage[]{microtype}
  \UseMicrotypeSet[protrusion]{basicmath} % disable protrusion for tt fonts
}{}
\makeatletter
\@ifundefined{KOMAClassName}{% if non-KOMA class
  \IfFileExists{parskip.sty}{%
    \usepackage{parskip}
  }{% else
    \setlength{\parindent}{0pt}
    \setlength{\parskip}{6pt plus 2pt minus 1pt}}
}{% if KOMA class
  \KOMAoptions{parskip=half}}
\makeatother
\usepackage{xcolor}
\usepackage[margin=20mm]{geometry}
\usepackage{color}
\usepackage{fancyvrb}
\newcommand{\VerbBar}{|}
\newcommand{\VERB}{\Verb[commandchars=\\\{\}]}
\DefineVerbatimEnvironment{Highlighting}{Verbatim}{commandchars=\\\{\}}
% Add ',fontsize=\small' for more characters per line
\newenvironment{Shaded}{}{}
\newcommand{\AlertTok}[1]{\textcolor[rgb]{1.00,0.00,0.00}{\textbf{#1}}}
\newcommand{\AnnotationTok}[1]{\textcolor[rgb]{0.38,0.63,0.69}{\textbf{\textit{#1}}}}
\newcommand{\AttributeTok}[1]{\textcolor[rgb]{0.49,0.56,0.16}{#1}}
\newcommand{\BaseNTok}[1]{\textcolor[rgb]{0.25,0.63,0.44}{#1}}
\newcommand{\BuiltInTok}[1]{\textcolor[rgb]{0.00,0.50,0.00}{#1}}
\newcommand{\CharTok}[1]{\textcolor[rgb]{0.25,0.44,0.63}{#1}}
\newcommand{\CommentTok}[1]{\textcolor[rgb]{0.38,0.63,0.69}{\textit{#1}}}
\newcommand{\CommentVarTok}[1]{\textcolor[rgb]{0.38,0.63,0.69}{\textbf{\textit{#1}}}}
\newcommand{\ConstantTok}[1]{\textcolor[rgb]{0.53,0.00,0.00}{#1}}
\newcommand{\ControlFlowTok}[1]{\textcolor[rgb]{0.00,0.44,0.13}{\textbf{#1}}}
\newcommand{\DataTypeTok}[1]{\textcolor[rgb]{0.56,0.13,0.00}{#1}}
\newcommand{\DecValTok}[1]{\textcolor[rgb]{0.25,0.63,0.44}{#1}}
\newcommand{\DocumentationTok}[1]{\textcolor[rgb]{0.73,0.13,0.13}{\textit{#1}}}
\newcommand{\ErrorTok}[1]{\textcolor[rgb]{1.00,0.00,0.00}{\textbf{#1}}}
\newcommand{\ExtensionTok}[1]{#1}
\newcommand{\FloatTok}[1]{\textcolor[rgb]{0.25,0.63,0.44}{#1}}
\newcommand{\FunctionTok}[1]{\textcolor[rgb]{0.02,0.16,0.49}{#1}}
\newcommand{\ImportTok}[1]{\textcolor[rgb]{0.00,0.50,0.00}{\textbf{#1}}}
\newcommand{\InformationTok}[1]{\textcolor[rgb]{0.38,0.63,0.69}{\textbf{\textit{#1}}}}
\newcommand{\KeywordTok}[1]{\textcolor[rgb]{0.00,0.44,0.13}{\textbf{#1}}}
\newcommand{\NormalTok}[1]{#1}
\newcommand{\OperatorTok}[1]{\textcolor[rgb]{0.40,0.40,0.40}{#1}}
\newcommand{\OtherTok}[1]{\textcolor[rgb]{0.00,0.44,0.13}{#1}}
\newcommand{\PreprocessorTok}[1]{\textcolor[rgb]{0.74,0.48,0.00}{#1}}
\newcommand{\RegionMarkerTok}[1]{#1}
\newcommand{\SpecialCharTok}[1]{\textcolor[rgb]{0.25,0.44,0.63}{#1}}
\newcommand{\SpecialStringTok}[1]{\textcolor[rgb]{0.73,0.40,0.53}{#1}}
\newcommand{\StringTok}[1]{\textcolor[rgb]{0.25,0.44,0.63}{#1}}
\newcommand{\VariableTok}[1]{\textcolor[rgb]{0.10,0.09,0.49}{#1}}
\newcommand{\VerbatimStringTok}[1]{\textcolor[rgb]{0.25,0.44,0.63}{#1}}
\newcommand{\WarningTok}[1]{\textcolor[rgb]{0.38,0.63,0.69}{\textbf{\textit{#1}}}}
\usepackage{longtable,booktabs,array}
\usepackage{calc} % for calculating minipage widths
% Correct order of tables after \paragraph or \subparagraph
\usepackage{etoolbox}
\makeatletter
\patchcmd\longtable{\par}{\if@noskipsec\mbox{}\fi\par}{}{}
\makeatother
% Allow footnotes in longtable head/foot
\IfFileExists{footnotehyper.sty}{\usepackage{footnotehyper}}{\usepackage{footnote}}
\makesavenoteenv{longtable}
\setlength{\emergencystretch}{3em} % prevent overfull lines
\providecommand{\tightlist}{%
  \setlength{\itemsep}{0pt}\setlength{\parskip}{0pt}}
\setcounter{secnumdepth}{-\maxdimen} % remove section numbering
% Markdown-inspired typography; compile with pdfLaTeX from this directory.
\usepackage{graphicx}
\usepackage{tgheros}
\usepackage[scale=1.0]{inconsolata}
\renewcommand{\familydefault}{\sfdefault}
\usepackage{seqsplit}
\DeclareRobustCommand{\InlineCode}[1]{{\ttfamily\seqsplit{#1}}}
\usepackage{titlesec}
\usepackage{needspace}
\definecolor{MarkdownText}{HTML}{24292F}
\definecolor{MarkdownRule}{HTML}{D0D7DE}
\definecolor{shadecolor}{HTML}{F6F8FA}
\color{MarkdownText}
\titleformat{\section}{\LARGE\bfseries\raggedright}{}{0pt}{}[\color{MarkdownRule}\titlerule]
\titleformat{\subsection}{\Large\bfseries\raggedright}{}{0pt}{}[\color{MarkdownRule}\titlerule]
\titleformat{\subsubsection}{\large\bfseries\raggedright}{}{0pt}{}
\titlespacing*{\section}{0pt}{1.3em}{0.7em}
\titlespacing*{\subsection}{0pt}{1.2em}{0.6em}
\titlespacing*{\subsubsection}{0pt}{1em}{0.4em}
\setlength{\parindent}{0pt}
\setlength{\parskip}{0.55em plus 0.1em minus 0.1em}
\renewcommand{\arraystretch}{1.18}
\setlength{\tabcolsep}{5pt}
\setlength{\emergencystretch}{3em}
\raggedright
\raggedbottom
\widowpenalty=10000
\clubpenalty=10000
\usepackage{framed}
\renewenvironment{Shaded}{\begin{snugshade}}{\end{snugshade}}
\usepackage{fvextra}
\fvset{breaklines=true,breakanywhere=true,fontsize=\small}
\DefineVerbatimEnvironment{verbatim}{Verbatim}{breaklines=true,breakanywhere=true,fontsize=\small}
\BeforeBeginEnvironment{verbatim}{\begin{snugshade}}
\AfterEndEnvironment{verbatim}{\end{snugshade}}
\AtBeginEnvironment{longtable}{\small}
\AtBeginDocument{\pdfstringdefDisableCommands{\def\InlineCode#1{#1}}}
\ifLuaTeX
  \usepackage{selnolig}  % disable illegal ligatures
\fi
\IfFileExists{bookmark.sty}{\usepackage{bookmark}}{\usepackage{hyperref}}
\IfFileExists{xurl.sty}{\usepackage{xurl}}{} % add URL line breaks if available
\urlstyle{same}
\hypersetup{
  colorlinks=true,
  linkcolor={blue},
  filecolor={Maroon},
  citecolor={Blue},
  urlcolor={blue},
  pdfcreator={LaTeX via pandoc}}

\author{}
\date{}

\begin{document}

\setstretch{1.06}
\hypertarget{understanding-loadbalancer.hpp-and-loadbalancer.cpp}{%
\section{\texorpdfstring{Understanding \InlineCode{LoadBalancer.hpp} and
\InlineCode{LoadBalancer.cpp}}{Understanding LoadBalancer.hpp and
LoadBalancer.cpp}}\label{understanding-loadbalancer.hpp-and-loadbalancer.cpp}}

This guide describes the implementation in
\href{LoadBalancer.hpp}{\InlineCode{LoadBalancer.hpp}} and
\href{LoadBalancer.cpp}{\InlineCode{LoadBalancer.cpp}} in
\InlineCode{sandbox/C++/hpc}, including its PBS extension. A separate
copy lives in \href{../../../umbridge/hpc}{\InlineCode{umbridge/hpc}};
the two copies are not identical. The explanations below follow the
local sandbox code and the bundled
\href{../../../umbridge/lib/umbridge.h}{\InlineCode{umbridge.h}}.

\hypertarget{purpose-and-overall-design}{%
\subsection{Purpose and overall
design}\label{purpose-and-overall-design}}

The load balancer exposes a UM-Bridge HTTP server to a client, such as a
numerical method requesting model evaluations. It obtains model servers
through an HPC scheduler and forwards operations to those servers as a
UM-Bridge client. This is the scope described in the
\href{https://um-bridge-benchmarks.readthedocs.io/en/docs/umbridge/hpc.html}{UM-Bridge
HPC overview}.

Startup discovers model names with a bootstrap job. If the first model
request uses the default \InlineCode{job.sh}, that job is adopted rather
than cancelled and replaced. Without a client session, later
\InlineCode{LoadBalancer} method calls still start separate jobs.
Clients can instead send a \InlineCode{UMBridge-Session} header; calls
for the same session and model reuse one server until
\InlineCode{POST /CloseSession}. The bundled C++ example uses this mode
so metadata, evaluations, and an optional Jacobian share one backend
job. A model-specific \InlineCode{job\_<name>.sh} takes precedence: the
bootstrap job is cancelled and that script is submitted instead. With
HyperQueue, jobs can run inside existing HPC allocations; one
HyperQueue job need not correspond to a new batch allocation.

The main division of responsibilities is:

\begin{longtable}[]{@{}
  >{\raggedright\arraybackslash}p{(\columnwidth - 2\tabcolsep) * \real{0.4300}}
  >{\raggedright\arraybackslash}p{(\columnwidth - 2\tabcolsep) * \real{0.5700}}@{}}
\toprule\noalign{}
\begin{minipage}[b]{\linewidth}\raggedright
Component
\end{minipage} & \begin{minipage}[b]{\linewidth}\raggedright
Responsibility
\end{minipage} \\
\midrule\noalign{}
\endhead
\bottomrule\noalign{}
\endlastfoot
\InlineCode{LoadBalancer} & Expose the UM-Bridge model API and acquire a
backend per call or reuse one by session. \\
\InlineCode{JobManager} / \InlineCode{CommandJobManager} & Combine
script selection, job submission, URL discovery, HTTP connection
creation, and session ownership. \\
\InlineCode{JobSubmitter} and its implementations & Translate
environment variables into scheduler arguments and create a
\InlineCode{Job}. \\
\InlineCode{Job} and its implementations & Hold a scheduler job
identifier and cancel the job on destruction. \\
\InlineCode{JobCommunicator} / \InlineCode{FilesystemCommunicator} &
Tell a job where to publish its server URL and read it back. \\
\InlineCode{JobCommunicatorFactory} /
\InlineCode{FilesystemCommunicatorFactory} & Create a separate
communicator for each submission. \\
\InlineCode{JobScriptLocator} & Select a default or model-specific
launch script. \\
\InlineCode{JobModel} & Keep the scheduler job alive for as long as its
HTTP model is owned by an operation or session. \\
\InlineCode{main} & Assemble the components, discover model names, and
start the public server. \\
\end{longtable}

Most of the implementation is in the header. The \InlineCode{.cpp} file
contains startup helpers and \InlineCode{main}; the split is
architectural rather than a conventional split between declarations and
definitions.

The included \InlineCode{umbridge.h} supplies the model interface, HTTP
client/server implementation, and \InlineCode{json} alias for
\InlineCode{nlohmann::json}. Standard-library types provide paths,
containers, streams, ownership pointers, synchronization, and timing.
Some of those declarations arrive through transitive includes rather
than explicit includes in \InlineCode{LoadBalancer.hpp}.

\hypertarget{a-request-from-start-to-finish}{%
\subsection{A request from start to
finish}\label{a-request-from-start-to-finish}}

For a sessioned client workflow, the central flow is:

\begin{verbatim}
Client -> LoadBalancer: ModelInfo / sizes / evaluations / optional Jacobian + session
LoadBalancer -> CommandJobManager: requestModelAccess(name, session_id)
Manager -> Scheduler: Submit selected script with URL directory
Scheduler -> Model server: Start job and model server
Model server -> Shared filesystem: Job script writes url-ID.txt
Manager -> Shared filesystem: Poll, then read URL
Manager -> Model server: HTTPModel discovers model and capabilities
Manager -> Shared filesystem: Communicator destruction removes URL file
Manager -> LoadBalancer: Session-owned JobModel
LoadBalancer -> Model server: Forward Evaluate through JobModel / HTTPModel
Model server -> LoadBalancer: Model outputs
LoadBalancer -> Client: Outputs
Client -> LoadBalancer: POST /CloseSession
LoadBalancer -> Manager: releaseSession(name, session_id)
Manager -> Scheduler: Job destruction requests cancellation
\end{verbatim}

The scheduler runs the script; the script, rather than the C++ load
balancer, chooses a port, starts the model executable, and publishes the
URL. The sequence groups those script actions with the backend for
readability.

In more detail, model discovery starts a bootstrap job. The first
request using the default script adopts it; subsequent fresh sessions
submit their own jobs.
\InlineCode{CommandJobManager::requestModelAccess(name,\space{}session\_id)}
then performs these steps:

\begin{enumerate}
\def\labelenumi{\arabic{enumi}.}
\tightlist
\item
  \InlineCode{locator.selectJobScript(name)} chooses a script. If it is
  the default \InlineCode{job.sh}, the request adopts the bootstrap job
  when available.
\item
  Otherwise, \InlineCode{job\_comm\_factory->create()} makes a
  communicator and \InlineCode{comm->getInitMessage()} supplies the
  environment to \InlineCode{job\_submitter->submit(...)}.
\item
  \InlineCode{comm->getModelUrl(job->getJobId())} waits for a newly
  submitted job's URL.
\item
  \InlineCode{umbridge::HTTPModel(url,\space{}name)} establishes the
  client-side model. Its constructor queries \InlineCode{/Info} to check
  the model name and \InlineCode{/ModelInfo} to read capabilities;
  constructing it already performs HTTP requests.
\item
  The job and HTTP model are moved into a \InlineCode{JobModel} retained
  by the session manager and shared with the caller. Later calls with
  the same session ID and model reuse this proxy rather than submitting
  another job.
\end{enumerate}

The bootstrap job is retained until adopted, the LoadBalancer exits, or
the scheduler's walltime expires. Start the client promptly after the
LoadBalancer if you want the first client session to use that job.

The client closes the session with \InlineCode{POST /CloseSession},
including the model name in the JSON body and the same session header.
The bundled \InlineCode{HTTPModel} accepts an optional session ID;
without one, existing clients retain the one-job-per-method behavior.
Session IDs are 1--128 ASCII letters, digits, underscores, or hyphens.
This session header and close endpoint are extensions implemented by
this local server/client pair, not part of the base UM-Bridge protocol.
The session is created lazily on the first per-model request; the
LoadBalancer's \InlineCode{/Info} endpoint returns the startup model-name
snapshot.

The local communicator is destroyed before
\InlineCode{requestModelAccess} returns, so its URL file is removed
\textbf{before} the forwarded evaluation begins. The URL has already
been consumed and the HTTP model retains the connection information. The
filesystem is used only to discover that URL: numerical inputs and
outputs travel over HTTP. The \InlineCode{HTTPModel} constructor uses
its default \InlineCode{useShMem\space{}=\space{}false}, so this backend
connection does not enable UM-Bridge's optional shared-memory transport.

Each overridden method holds a shared pointer to the returned model
while it forwards the operation. For sessioned requests the manager
retains another owner until \InlineCode{/CloseSession}; for requests
without a session, the local owner is the only one and its destruction
requests scheduler cancellation.

Without a session, the unit of job creation is a
\textbf{\InlineCode{LoadBalancer} method call}, not necessarily an
incoming HTTP request. Size and capability queries also create jobs.
The bundled HTTP \InlineCode{/ModelInfo} handler calls all four
\InlineCode{Supports*} methods. With a session, those calls and
subsequent size/evaluation requests all use the same backend. The
\InlineCode{/ApplyHessian} handler additionally calls
\InlineCode{GetInputSizes} to check the result size, even when
\InlineCode{error\_checks} is false.

\hypertarget{startup-and-command-line-handling-in-loadbalancer.cpp}{%
\subsection{\texorpdfstring{Startup and command-line handling in
\InlineCode{LoadBalancer.cpp}}{Startup and command-line handling in
LoadBalancer.cpp}}\label{startup-and-command-line-handling-in-loadbalancer.cpp}}

\hypertarget{clear_urldirectory}{%
\subsubsection{\texorpdfstring{\InlineCode{clear\_url(directory)}}{clear\_url(directory)}}\label{clear_urldirectory}}

This helper returns immediately if the directory is absent. Otherwise,
it visits its immediate entries and removes names matching
\InlineCode{url-\textbackslash{}d+\textbackslash{}.txt}: for example,
\InlineCode{url-123.txt}. It does not recurse, and a name containing a
scheduler suffix, such as \InlineCode{url-123.server.txt}, does not
match.

\InlineCode{main} calls it on \InlineCode{urls} before validating
command-line arguments. Its purpose is to prevent old URL files from
being mistaken for newly submitted jobs. The directory therefore must
belong to this load-balancer instance; sharing it with another active
instance allows one startup to remove another instance's files.
Filesystem failures are not caught here.

\hypertarget{launch_hq_with_alloc_queue}{%
\subsubsection{\texorpdfstring{\InlineCode{launch\_hq\_with\_alloc\_queue()}}{launch\_hq\_with\_alloc\_queue()}}\label{launch_hq_with_alloc_queue}}

For HyperQueue, startup runs shell commands to stop a local HQ server,
start one in the background, repeatedly query
\InlineCode{./hq\space{}server\space{}info} until it responds, and
execute \InlineCode{hq\_scripts/allocation\_queue.sh}.

The binary location \InlineCode{./hq} and the script path are relative
to the process's working directory. The commands use \InlineCode{\&>}
redirection, so they assume a shell supporting that syntax;
\InlineCode{std::system} does not itself guarantee Bash. Exit statuses
are ignored and the readiness loop has no timeout. There is no matching
HQ-server shutdown routine when \InlineCode{main} ends.

\hypertarget{get_argargs-arg_name}{%
\subsubsection{\texorpdfstring{\InlineCode{get\_arg(args,\space{}arg\_name)}}{\InlineCode{get\_arg(args,\space{}arg\_name)}}}\label{get_argargs-arg_name}}

The helper searches for the first string beginning with
\InlineCode{--<arg\_name>=} and having at least one character after
\InlineCode{=}. It returns that suffix, or an empty string when no
matching nonempty argument exists. Thus \InlineCode{--port=5000} is
accepted, while \InlineCode{--port\space{}5000} and \InlineCode{--port=}
are not treated as supplied values. Unrecognized arguments are not
rejected.

\begin{longtable}[]{@{}
  >{\raggedright\arraybackslash}p{(\columnwidth - 4\tabcolsep) * \real{0.4500}}
  >{\raggedright\arraybackslash}p{(\columnwidth - 4\tabcolsep) * \real{0.1400}}
  >{\raggedright\arraybackslash}p{(\columnwidth - 4\tabcolsep) * \real{0.4100}}@{}}
\toprule\noalign{}
\begin{minipage}[b]{\linewidth}\raggedright
Option
\end{minipage} & \begin{minipage}[b]{\linewidth}\raggedright
Default
\end{minipage} & \begin{minipage}[b]{\linewidth}\raggedright
Use
\end{minipage} \\
\midrule\noalign{}
\endhead
\bottomrule\noalign{}
\endlastfoot
\InlineCode{--scheduler=hyperqueue}, \InlineCode{--scheduler=slurm}, or
\InlineCode{--scheduler=pbs} & Required & Select the submitter and
script directory. See the PBS naming issue below. \\
\InlineCode{--delay-ms=N} & \InlineCode{0} & Set the optional delay used
by the submitter. \\
\InlineCode{--port=N} & \InlineCode{4242} & Select the public
load-balancer TCP port. \\
\end{longtable}

The two numeric options use \InlineCode{std::stoi}. Invalid or
out-of-range conversions can throw uncaught exceptions; the code has no
explicit range checks and does not verify that the whole string was
consumed. A negative submission delay does not enter the
\InlineCode{delay\space{}>\space{}0} branch.

\hypertarget{mainargc-argv}{%
\subsubsection{\texorpdfstring{\InlineCode{main(argc,\space{}argv)}}{\InlineCode{main(argc,\space{}argv)}}}\label{mainargc-argv}}

After cleaning the URL directory, \InlineCode{main} reads the options
and constructs the selected scheduler components:

\begin{longtable}[]{@{}
  >{\raggedright\arraybackslash}p{(\columnwidth - 4\tabcolsep) * \real{0.2200}}
  >{\raggedright\arraybackslash}p{(\columnwidth - 4\tabcolsep) * \real{0.5000}}
  >{\raggedright\arraybackslash}p{(\columnwidth - 4\tabcolsep) * \real{0.2800}}@{}}
\toprule\noalign{}
\begin{minipage}[b]{\linewidth}\raggedright
Scheduler value
\end{minipage} & \begin{minipage}[b]{\linewidth}\raggedright
Submitter class intended by the design
\end{minipage} & \begin{minipage}[b]{\linewidth}\raggedright
Script directory
\end{minipage} \\
\midrule\noalign{}
\endhead
\bottomrule\noalign{}
\endlastfoot
\InlineCode{hyperqueue} & \InlineCode{HyperQueueSubmitter} &
\InlineCode{hq\_scripts} \\
\InlineCode{slurm} & \InlineCode{SlurmSubmitter} &
\InlineCode{slurm\_scripts} \\
\InlineCode{pbs} & \InlineCode{PBSSubmitter} &
\InlineCode{pbs\_scripts} \\
\end{longtable}

Missing or unsupported scheduler values print an error and call
\InlineCode{std::exit(-1)}. The current PBS branch spells the class
\InlineCode{PbsSubmitter}, whereas the header declares
\InlineCode{PBSSubmitter}; this is a compile-time error in the supplied
source.

Startup then:

\begin{enumerate}
\def\labelenumi{\arabic{enumi}.}
\tightlist
\item
  Creates a \InlineCode{FilesystemCommunicatorFactory} for
  \InlineCode{urls}, polling every 500 ms. Its constructor creates the
  directory if needed.
\item
  Configures \InlineCode{JobScriptLocator} with default
  \InlineCode{job.sh} and model-specific names
  \InlineCode{job\_<model\_name>.sh}.
\item
  Moves the submitter and communicator factory into a shared
  \InlineCode{CommandJobManager}.
\item
  Calls \InlineCode{getModelNames()}. This submits a real job using the
  default script, waits for its server, and calls
  \InlineCode{umbridge::SupportedModels(url)}. The manager retains the
  job as a bootstrap backend. The public server is not listening while
  this discovery job is pending.
\item
  Prints each model's selected script and creates one
  \InlineCode{LoadBalancer} object per discovered model name. The
  printed list is sorted using a copy of the names; the original
  discovery order is retained in the model vector.
\item
  Builds a vector of non-owning \InlineCode{umbridge::Model*} pointers
  into the completed \InlineCode{LB\_vector}, then calls
  \InlineCode{umbridge::serveModels}.
\end{enumerate}

The pointer conversion adapts to the UM-Bridge server API.
\InlineCode{LB\_vector} remains alive for the blocking server call, and
no later insertion invalidates its element addresses.

The final call is:

\begin{Shaded}
\begin{Highlighting}[]
\NormalTok{umbridge}\OperatorTok{::}\NormalTok{serveModels}\OperatorTok{(}\NormalTok{LB\_ptr\_vector}\OperatorTok{,} \StringTok{"0.0.0.0"}\OperatorTok{,}\NormalTok{ port}\OperatorTok{,} \KeywordTok{true}\OperatorTok{,} \KeywordTok{false}\OperatorTok{);}
\end{Highlighting}
\end{Shaded}

\InlineCode{0.0.0.0} binds on all IPv4 interfaces. \InlineCode{true}
sets \InlineCode{enable\_parallel}, allowing HTTP handlers to call
models without the library's serializing model lock. \InlineCode{false}
sets \InlineCode{error\_checks}, disabling the checks guarded by that
flag. Those checks would otherwise call size and capability methods;
without a session, those extra calls can submit additional jobs. It does
not disable every check in the bundled library, and the load-balancer
classes themselves do not validate numerical inputs.

\hypertarget{utilities-and-command-construction-in-loadbalancer.hpp}{%
\subsection{\texorpdfstring{Utilities and command construction in
\InlineCode{LoadBalancer.hpp}}{Utilities and command construction in
LoadBalancer.hpp}}\label{utilities-and-command-construction-in-loadbalancer.hpp}}

\hypertarget{get_command_outputcommand}{%
\subsubsection{\texorpdfstring{\InlineCode{get\_command\_output(command)}}{get\_command\_output(command)}}\label{get_command_outputcommand}}

\InlineCode{popen(command.c\_str(),\space{}"r")} runs a shell command
and opens its standard output for reading. A
\InlineCode{std::unique\_ptr<FILE,\space{}decltype(\&pclose)>} couples
the pipe to the \InlineCode{pclose} cleanup function. \InlineCode{fgets}
reads chunks into a 128-byte buffer until no more data is available, and
the chunks are concatenated into the returned string. The buffer limits
each read, not the total output size.

If \InlineCode{popen} cannot create the pipe, the helper throws
\InlineCode{std::runtime\_error} with the command and \InlineCode{errno}
description. It does not capture standard error unless the command
redirects it, and it discards the command's exit status returned by
\InlineCode{pclose}. It also does not inspect the stream error indicator
when reading stops. A scheduler submission failure can therefore produce
empty or unexpected output without an exception from this helper.

\hypertarget{other-small-helpers}{%
\subsubsection{Other small helpers}\label{other-small-helpers}}

\begin{longtable}[]{@{}
  >{\raggedright\arraybackslash}p{(\columnwidth - 2\tabcolsep) * \real{0.4700}}
  >{\raggedright\arraybackslash}p{(\columnwidth - 2\tabcolsep) * \real{0.5300}}@{}}
\toprule\noalign{}
\begin{minipage}[b]{\linewidth}\raggedright
Function
\end{minipage} & \begin{minipage}[b]{\linewidth}\raggedright
Exact behavior
\end{minipage} \\
\midrule\noalign{}
\endhead
\bottomrule\noalign{}
\endlastfoot
\InlineCode{wait\_for\_file(file\_path,\space{}polling\_cycle)} & Check
\InlineCode{std::filesystem::exists} repeatedly and sleep between failed
checks. It has no timeout or scheduler-state check. \\
\InlineCode{read\_line\_from\_file(file\_path)} & Open the file, throw
if opening fails, then return its first line. It does not verify
nonempty contents or successful line extraction. \\
\InlineCode{remove\_trailing\_newline(s)} & Remove at most one final
\InlineCode{\textbackslash{}n}; it does not trim spaces or a preceding
\InlineCode{\textbackslash{}r}. \\
\end{longtable}

\hypertarget{command}{%
\subsubsection{\texorpdfstring{\InlineCode{Command}}{Command}}\label{command}}

This aggregate holds \InlineCode{exec}, a vector of
\InlineCode{options}, and \InlineCode{target}. \InlineCode{addOption}
appends an option. \InlineCode{toString} joins the executable, each
option, and the target with spaces. For example, a \InlineCode{SlurmJob}
constructs a command with \InlineCode{exec\space{}=\space{}"sbatch"},
environment options, \InlineCode{--parsable}, and the script path.

These are shell command fragments, not an argument vector passed
directly to an executable. \InlineCode{Command} performs no quoting or
escaping. Script paths and environment values containing shell syntax
require care, and untrusted input must not be interpolated into this
command path. The PBS environment helper provides shell quoting for its
environment list only.

\hypertarget{scheduler-jobs-and-submitters}{%
\subsection{Scheduler jobs and
submitters}\label{scheduler-jobs-and-submitters}}

\hypertarget{job-resource-ownership}{%
\subsubsection{\texorpdfstring{\InlineCode{Job}: resource
ownership}{Job: resource
ownership}}\label{job-resource-ownership}}

\InlineCode{Job} is an abstract owner of a scheduler job, exposing
\InlineCode{getJobId()}. Its virtual destructor allows correct
destruction through a base pointer. Copy construction, move
construction, and both assignments are deleted so two \InlineCode{Job}
objects cannot accidentally claim the same cancellation responsibility.

Ownership is transferred by moving a \InlineCode{std::unique\_ptr<Job>},
which moves the pointer rather than the \InlineCode{Job} itself. The
concrete constructors submit jobs; their destructors invoke scheduler
cancellation commands. This is RAII (Resource Acquisition Is
Initialization): a resource's lifetime follows the lifetime of its
owning C++ object.

\begin{longtable}[]{@{}
  >{\raggedright\arraybackslash}p{(\columnwidth - 6\tabcolsep) * \real{0.1900}}
  >{\raggedright\arraybackslash}p{(\columnwidth - 6\tabcolsep) * \real{0.3200}}
  >{\raggedright\arraybackslash}p{(\columnwidth - 6\tabcolsep) * \real{0.2900}}
  >{\raggedright\arraybackslash}p{(\columnwidth - 6\tabcolsep) * \real{0.2000}}@{}}
\toprule\noalign{}
\begin{minipage}[b]{\linewidth}\raggedright
Class
\end{minipage} & \begin{minipage}[b]{\linewidth}\raggedright
Submission
\end{minipage} & \begin{minipage}[b]{\linewidth}\raggedright
ID extraction
\end{minipage} & \begin{minipage}[b]{\linewidth}\raggedright
Destruction
\end{minipage} \\
\midrule\noalign{}
\endhead
\bottomrule\noalign{}
\endlastfoot
\InlineCode{HyperQueueJob} &
\InlineCode{./hq\space{}submit\space{}...\space{}--output-mode=quiet\space{}<script>}
& Keep stdout after removing one trailing newline. &
\InlineCode{./hq\space{}job\space{}cancel\space{}<id>} \\
\InlineCode{SlurmJob} &
\InlineCode{sbatch\space{}...\space{}--parsable\space{}<script>} & Match
a numeric ID, optionally followed by \InlineCode{;<cluster>}, and keep
the digits. & \InlineCode{scancel\space{}<id>} \\
\InlineCode{PBSJob} & \InlineCode{qsub\space{}...\space{}<script>} &
Match a numeric ID, optionally followed by \InlineCode{.<server>}, and
keep the digits. & \InlineCode{qdel\space{}<id>} \\
\end{longtable}

The SLURM/PBS parsers examine output line by line, ignore nonmatching
lines, and retain the last matching ID. Their optional suffix permits
only letters, digits, \InlineCode{\_}, and \InlineCode{-}; PBS array IDs
or suffixes containing additional dots do not match. No matching line
leaves an empty ID. Discarding the suffix also means subsequent
cancellation and URL lookup use only the numeric component.

Each \InlineCode{getJobId()} simply returns the stored string. The
destructors ignore cancellation exit codes: destruction requests
cancellation but does not verify that the remote resources were
released.

\hypertarget{jobsubmitter-scheduling-policy-and-environment-conversion}{%
\subsubsection{\texorpdfstring{\InlineCode{JobSubmitter}: scheduling
policy and environment
conversion}{JobSubmitter: scheduling
policy and environment
conversion}}\label{jobsubmitter-scheduling-policy-and-environment-conversion}}

\InlineCode{JobSubmitter::submit(job\_script,\space{}env)} takes a
script and a map of environment variables and returns exclusive
ownership of the resulting \InlineCode{Job}. Its virtual destructor
supports deletion through the interface.

All three implementations store a configured delay and a mutex. For a
positive delay, \InlineCode{submit} locks the mutex, sleeps, and then
leaves that block. \textbf{The mutex is released before constructing the
job and executing the submission command.} It serializes the waiting
sections, not entire submissions, and does not enforce a strict minimum
interval between actual scheduler command invocations.

The private \InlineCode{env\_to\_options} methods generate different
command fragments:

\begin{longtable}[]{@{}
  >{\raggedright\arraybackslash}p{(\columnwidth - 4\tabcolsep) * \real{0.2600}}
  >{\raggedright\arraybackslash}p{(\columnwidth - 4\tabcolsep) * \real{0.3500}}
  >{\raggedright\arraybackslash}p{(\columnwidth - 4\tabcolsep) * \real{0.3900}}@{}}
\toprule\noalign{}
\begin{minipage}[b]{\linewidth}\raggedright
Submitter
\end{minipage} & \begin{minipage}[b]{\linewidth}\raggedright
Environment encoding
\end{minipage} & \begin{minipage}[b]{\linewidth}\raggedright
Additional behavior
\end{minipage} \\
\midrule\noalign{}
\endhead
\bottomrule\noalign{}
\endlastfoot
\InlineCode{HyperQueueSubmitter} & One
\InlineCode{--env\space{}KEY=VALUE} option per map entry. & Add
\InlineCode{--priority=-<job\_count>}; increment the count after
creating the job. \\
\InlineCode{SlurmSubmitter} & One
\InlineCode{--export=ALL,KEY=VALUE,...} option. & Include the inherited
environment via \InlineCode{ALL}. \\
\InlineCode{PBSSubmitter} & \InlineCode{-v} followed by a single-quoted
\InlineCode{KEY=VALUE,...} list. & Return no options for an empty map;
escape embedded single quotes for the shell. \\
\end{longtable}

The HyperQueue priority becomes increasingly negative as the count
increases, expressing a preference for earlier submissions.
\InlineCode{job\_count} is atomic, so reads and increments do not race
at the C++ memory-access level. Reading its value and later incrementing
it are separate operations, however, so concurrent calls can generate
the same priority. It is not a unique job identifier.

PBS shell quoting preserves the list as one shell argument; it does not
escape the comma separators of PBS's own environment-list syntax. The
generic command builder still leaves the script path unquoted.

\hypertarget{communicating-a-backend-url}{%
\subsection{Communicating a backend
URL}\label{communicating-a-backend-url}}

\hypertarget{the-interfaces}{%
\subsubsection{The interfaces}\label{the-interfaces}}

\InlineCode{JobCommunicator} defines two phases:

\begin{itemize}
\tightlist
\item
  \InlineCode{getInitMessage()} returns key-value pairs supplied to the
  job environment.
\item
  \InlineCode{getModelUrl(job\_id)} obtains the URL published by that
  job.
\end{itemize}

Its copy/move operations are deleted for the same ownership reason as
\InlineCode{Job}: each communicator owns its cleanup responsibility.
\InlineCode{JobCommunicatorFactory} exposes \InlineCode{create()} to
obtain a fresh owning pointer for every job, avoiding shared mutable
per-job state.

\hypertarget{filesystemcommunicator}{%
\subsubsection{\texorpdfstring{\InlineCode{FilesystemCommunicator}}{FilesystemCommunicator}}\label{filesystemcommunicator}}

The constructor stores \InlineCode{file\_dir} and the polling interval.
\InlineCode{getInitMessage()} returns a map containing:

\begin{Shaded}
\begin{Highlighting}[]
\NormalTok{UMBRIDGE\_LOADBALANCER\_COMM\_FILEDIR=\textless{}file\_dir\textgreater{}}
\end{Highlighting}
\end{Shaded}

\InlineCode{getUrlFileName(job\_id)} forms
\InlineCode{url-<job\_id>.txt}. \InlineCode{getModelUrl} stores the full
path, logs it, waits for existence, and returns the first line. The
destructor removes that path if one has been assigned. It does not
remove the directory.

\InlineCode{FilesystemCommunicatorFactory} stores the common directory
and polling interval, creates the directory in its constructor, and
returns a new communicator from \InlineCode{create()}. It calls
\InlineCode{create\_directory}, not recursive
\InlineCode{create\_directories}.

\hypertarget{contract-with-a-launch-script}{%
\subsubsection{Contract with a launch
script}\label{contract-with-a-launch-script}}

The C++ code assumes a script will:

\begin{enumerate}
\def\labelenumi{\arabic{enumi}.}
\tightlist
\item
  Start a UM-Bridge server in its allocation and keep the job alive
  while that server is needed.
\item
  Determine a URL reachable from the load-balancer host.
\item
  Publish that URL as the first line of
  \InlineCode{\$UMBRIDGE\_LOADBALANCER\_COMM\_FILEDIR/url-<id>.txt},
  using the same ID string returned by \InlineCode{Job::getJobId()}.
\item
  Keep the server running until the load balancer cancels the job.
\end{enumerate}

The reference \href{../../../umbridge/hpc/slurm_scripts/job.sh}{SLURM
script} and \href{../../../umbridge/hpc/hq_scripts/job.sh}{HyperQueue
script} select a free port, export \InlineCode{PORT}, start
\InlineCode{./testmodel}, wait for an HTTP response, write the URL using
\InlineCode{SLURM\_JOB\_ID} or \InlineCode{HQ\_JOB\_ID}, and finish with
\InlineCode{sleep\space{}infinity}. Hardware resource requests and the
model launch command are configured in those scripts, not in
\InlineCode{LoadBalancer}.

Both processes must resolve the URL directory to the same shared
filesystem location. Since \InlineCode{main} supplies the relative path
\InlineCode{urls}, the job's working directory matters. A PBS script
must normalize its job ID consistently with \InlineCode{PBSJob}, which
strips the server suffix. This sandbox directory currently has no
\InlineCode{pbs\_scripts/job.sh} implementing that contract.

The C++ existence check is not a complete readiness protocol. It does
not wait for a nonempty URL, validate its format, or retry a failed HTTP
connection. A file can become visible before the writer finishes it. A
script that publishes a completed temporary file by renaming it can
avoid that particular partial-file window, but the reference scripts
currently write directly to the final path.

\hypertarget{script-selection-and-model-management}{%
\subsection{Script selection and model
management}\label{script-selection-and-model-management}}

\hypertarget{jobscriptlocator}{%
\subsubsection{\texorpdfstring{\InlineCode{JobScriptLocator}}{JobScriptLocator}}\label{jobscriptlocator}}

The public fields describe the directory, default filename, model-name
prefix, and model-name suffix. With the values from \InlineCode{main},
the model \InlineCode{heat} selects
\InlineCode{<script\_dir>/job\_heat.sh} if it exists, otherwise
\InlineCode{<script\_dir>/job.sh}.

\InlineCode{selectJobScript(model\_name)} checks those paths in that
order and throws \InlineCode{std::runtime\_error} if neither exists. It
checks existence, not executable permission or script validity.
\InlineCode{getDefaultJobScript()} only constructs the default path; it
does not check existence.
\InlineCode{printModelJobScripts(model\_names)} sorts its by-value copy
and prints each selected path, potentially throwing if a required path
is absent.

The startup discovery job always uses the default script. Model-specific
scripts can give known models different resources, but cannot replace
the default discovery job. Model names are inserted directly into
filenames and are assumed to be suitable for this naming convention.

\hypertarget{jobmanager-and-commandjobmanager}{%
\subsubsection{\texorpdfstring{\InlineCode{JobManager} and
\InlineCode{CommandJobManager}}{JobManager and
CommandJobManager}}\label{jobmanager-and-commandjobmanager}}

\InlineCode{JobManager} abstracts model discovery, model acquisition,
and session release. \InlineCode{requestModelAccess(name,\space{}session\_id)}
returns a shared model; with an empty session ID its local owner
controls the job lifetime, while a nonempty ID causes the manager to
retain the backend until
\InlineCode{releaseSession(name,\space{}session\_id)}.

\InlineCode{CommandJobManager} owns one submitter, one communicator
factory, and a locator value. Its constructor accepts the first two
through \InlineCode{unique\_ptr} and moves them into members. The
acquisition procedure was described above.

\InlineCode{getModelNames()} uses the same submission/communication
mechanisms, but invokes \InlineCode{umbridge::SupportedModels(url)}
directly rather than constructing a \InlineCode{JobModel}. The function
retains the submitted job and URL as a bootstrap backend. The
communicator removes the URL file after the URL has been read; the server
stays alive. The first request using the default script adopts that job,
while a model-specific script cancels it and submits the selected script
instead.

Neither operation receives model inputs or \InlineCode{config\_json}.
Job scripts are selected by model name alone, so numerical configuration
cannot dynamically change resource requests through this interface.

\hypertarget{jobmodel-couple-a-model-to-its-job}{%
\subsubsection{\texorpdfstring{\InlineCode{JobModel}: couple a model to
its
job}{JobModel: couple a model to
its
job}}\label{jobmodel-couple-a-model-to-its-job}}

\InlineCode{JobModel} derives from \InlineCode{umbridge::Model} and
stores both a \InlineCode{unique\_ptr<Job>} and a
\InlineCode{unique\_ptr<umbridge::Model>}. Its constructor copies the
wrapped model's name into the base class, then takes ownership of both
pointers. Both are expected to refer to valid objects; the constructor
dereferences the model pointer.

All numerical, size, and capability methods delegate to the stored model
without changing arguments or results. This repetition preserves the
full virtual interface while adding lifetime management around it.

No explicit destructor is necessary: members are destroyed in reverse
declaration order. \InlineCode{model} is declared after
\InlineCode{job}, so the HTTP model is destroyed first and the scheduler
job second. A virtual destructor inherited from
\InlineCode{umbridge::Model} makes this work when the wrapper is held as
a base pointer.

\hypertarget{the-numerical-interface-exposed-by-loadbalancer}{%
\subsection{\texorpdfstring{The numerical interface exposed by
\InlineCode{LoadBalancer}}{The numerical interface exposed by
LoadBalancer}}\label{the-numerical-interface-exposed-by-loadbalancer}}

The constructor stores a model name in the \InlineCode{umbridge::Model}
base and shares the \InlineCode{JobManager} via
\InlineCode{std::shared\_ptr}. All named load balancers therefore use
the same submission infrastructure. The constructor itself does not
submit a job.

The numerical data use blocks:
\InlineCode{std::vector<std::vector<double>>} represents a list of input
or output vectors, not necessarily a rectangular matrix. For example,
input sizes \InlineCode{\{3,\space{}1\}} describe two input blocks with
lengths three and one. Indices such as \InlineCode{inWrt} and
\InlineCode{outWrt} select blocks, starting at zero.

Let output block \InlineCode{F\_i(x)} have length \InlineCode{m\_i}, and
input block \InlineCode{x\_j} length \InlineCode{n\_j}. Write its
Jacobian block as \InlineCode{J\_ij\space{}=\space{}dF\_i/dx\_j}.

\begin{longtable}[]{@{}
  >{\raggedright\arraybackslash}p{(\columnwidth - 2\tabcolsep) * \real{0.5100}}
  >{\raggedright\arraybackslash}p{(\columnwidth - 2\tabcolsep) * \real{0.4900}}@{}}
\toprule\noalign{}
\begin{minipage}[b]{\linewidth}\raggedright
Method
\end{minipage} & \begin{minipage}[b]{\linewidth}\raggedright
Meaning and returned data
\end{minipage} \\
\midrule\noalign{}
\endhead
\bottomrule\noalign{}
\endlastfoot
\InlineCode{GetInputSizes(config\_json)} & Length of every input block
under the supplied model configuration. \\
\InlineCode{GetOutputSizes(config\_json)} & Length of every output block
under that configuration. \\
\InlineCode{Evaluate(inputs,\space{}config\_json)} & Evaluate all output
blocks at the input blocks. \\
\InlineCode{Gradient(outWrt,\space{}inWrt,\space{}inputs,\space{}sens,\space{}config\_json)}
& Return \InlineCode{J\_ij\textasciicircum{}T\space{}sens}, an
\InlineCode{n\_j}-element vector. \InlineCode{sens} has
\InlineCode{m\_i} entries. \\
\InlineCode{ApplyJacobian(outWrt,\space{}inWrt,\space{}inputs,\space{}vec,\space{}config\_json)}
& Return \InlineCode{J\_ij\space{}vec}, an \InlineCode{m\_i}-element
vector. \InlineCode{vec} has \InlineCode{n\_j} entries. \\
\InlineCode{ApplyHessian(outWrt,\space{}inWrt1,\space{}inWrt2,\space{}inputs,\space{}sens,\space{}vec,\space{}config\_json)}
& Apply the mixed Hessian of the scalar
\InlineCode{sens\textasciicircum{}T\space{}F\_i(x)} to \InlineCode{vec}
in input block \InlineCode{inWrt2}; return a vector in block
\InlineCode{inWrt1}. \\
\InlineCode{SupportsEvaluate()} & Report whether evaluation is
supported. \\
\InlineCode{SupportsGradient()} & Report whether the gradient operation
is supported. \\
\InlineCode{SupportsApplyJacobian()} & Report whether a Jacobian action
is supported. \\
\InlineCode{SupportsApplyHessian()} & Report whether a Hessian action is
supported. \\
\end{longtable}

For the Hessian action, the relevant matrix has dimensions
\InlineCode{n\_inWrt1\space{}\texttimes{}\space{}n\_inWrt2}.
\InlineCode{sens} selects a weighted combination of entries in output
block \InlineCode{outWrt}; \InlineCode{vec} supplies a direction in
input block \InlineCode{inWrt2}.

\InlineCode{config\_json} is the model-specific JSON configuration,
defaulting to an empty object at these declarations. The load balancer
forwards it unchanged. It does not compute derivatives, approximate
unsupported operations, or interpret model parameters. Those
responsibilities belong to the backend.

\InlineCode{JobModel} forwards each method to its backend.
\InlineCode{LoadBalancer} obtains that backend through the manager.
Without a session, every size or support method still has scheduling
cost. With a session, those requests reuse the session's backend; the
\InlineCode{HTTPModel} at the client side also caches capability flags
after the first \InlineCode{/ModelInfo} request.

\hypertarget{concurrency-ownership-and-failures}{%
\subsection{Concurrency, ownership, and
failures}\label{concurrency-ownership-and-failures}}

The public server can process concurrent requests because
\InlineCode{enable\_parallel} is true. Different session IDs own
different backend jobs; requests carrying the same session ID reuse that
session's model proxy. A client should use a unique session ID and issue
calls for it sequentially. The HPC scheduler decides when resources
become available. The C++ code does not implement its own job queue,
admission limit, or load metric.

Shared state is limited mainly to the manager's components: a
submitter's mutex and, for HyperQueue, atomic counter. Each communicator
has its own \InlineCode{file\_path}. Correct filesystem isolation still
depends on IDs and paths being unique and consistent. The delay mutex is
not held while an evaluation runs.

Normal return and ordinary exception unwinding destroy owning pointers.
If URL reading or HTTP-model construction throws after a job was
acquired, its local owners request cancellation on the way out. A
session-owned job is cancelled when the client closes the session; the
example does this on normal completion and during exception unwinding.

There are limits to that guarantee:

\begin{itemize}
\tightlist
\item
  A failed submission can leave an empty or malformed ID because command
  exit codes and parsed IDs are not validated.
\item
  Waiting for a URL has no timeout, cancellation signal, or job-failure
  check; a job that never publishes a URL can leave the request waiting
  indefinitely.
\item
  A forced process termination does not run C++ destructors, and this
  code installs no signal-driven cleanup handler. A forcibly terminated
  client therefore cannot send \InlineCode{/CloseSession}; the
  scheduler's job walltime is the eventual backstop.
\item
  Scheduler cancellation failures are ignored. The filesystem destructor
  uses the throwing form of \InlineCode{std::filesystem::remove}; an
  exception escaping a destructor can terminate the process.
\item
  The code does not catch and translate its startup failures in
  \InlineCode{main}. Stack unwinding is not guaranteed when an exception
  escapes \InlineCode{main}, so cleanup cannot be assumed in that case.
  Runtime HTTP failure handling also depends on the bundled
  UM-Bridge/HTTP library.
\end{itemize}

These are descriptions of the present implementation, rather than
behavior added by the documentation changes.

\hypertarget{building-and-extending-this-local-copy}{%
\subsection{Building and extending this local
copy}\label{building-and-extending-this-local-copy}}

The code requires C++17 features, including
\InlineCode{std::filesystem}, structured bindings, and class-template
argument deduction. It relies on the bundled UM-Bridge header and its
\InlineCode{json.hpp}/\InlineCode{httplib.h} dependencies, threads,
POSIX process facilities, scheduler commands, and a shared URL
directory.

The sandbox copy is not currently a complete standalone deployment:

\begin{itemize}
\tightlist
\item
  \InlineCode{LoadBalancer.cpp} refers to \InlineCode{PbsSubmitter}, but
  the declared class is \InlineCode{PBSSubmitter}. C++ identifiers are
  case-sensitive, so this must be reconciled before the translation unit
  can compile.
\item
  The header includes \InlineCode{"umbridge.h"}, while the
  \InlineCode{.cpp} also includes \InlineCode{"../lib/umbridge.h"}.
  There is no \InlineCode{sandbox/C++/lib} directory. Suitable include
  search paths can resolve both to the bundled header:
  \InlineCode{-Iumbridge/lib} from the workspace root resolves the
  includes in the current layout. The relative include is not satisfied
  by the sandbox layout alone.
\item
  \href{../Makefile}{\InlineCode{sandbox/C++/Makefile}} builds
  \InlineCode{minimal-server} and \InlineCode{http-client}, not the load
  balancer. The separate \href{../../../umbridge/hpc/Makefile}{upstream
  HPC Makefile} provides a build recipe for its own directory layout.
\item
  \InlineCode{hq\_scripts}, \InlineCode{slurm\_scripts}, and
  \InlineCode{pbs\_scripts} are absent from this sandbox
  \InlineCode{hpc} directory. The first two have reference
  implementations in
  \href{../../../umbridge/hpc}{\InlineCode{umbridge/hpc}}; a compatible
  PBS model-job script is still needed for a PBS deployment.
\item
  \href{LoadBalancer.pbs}{\InlineCode{LoadBalancer.pbs}} is a batch
  launcher for the load-balancer executable, not the per-model
  URL-publishing script. Its current invocation also omits the required
  \InlineCode{--scheduler=...} option.
\end{itemize}

Once a working build and launch-script layout are supplied, the intended
command syntax is, for example:

\begin{Shaded}
\begin{Highlighting}[]
\ExtensionTok{./load{-}balancer} \AttributeTok{{-}{-}scheduler}\OperatorTok{=}\NormalTok{slurm }\AttributeTok{{-}{-}port}\OperatorTok{=}\NormalTok{4242 }\AttributeTok{{-}{-}delay{-}ms}\OperatorTok{=}\NormalTok{100}
\end{Highlighting}
\end{Shaded}

All relative paths are resolved against the launch directory, not
automatically against the source-file directory.

The interfaces expose several clear extension points:

\begin{longtable}[]{@{}
  >{\raggedright\arraybackslash}p{(\columnwidth - 2\tabcolsep) * \real{0.4000}}
  >{\raggedright\arraybackslash}p{(\columnwidth - 2\tabcolsep) * \real{0.6000}}@{}}
\toprule\noalign{}
\begin{minipage}[b]{\linewidth}\raggedright
Desired change
\end{minipage} & \begin{minipage}[b]{\linewidth}\raggedright
Main code to extend
\end{minipage} \\
\midrule\noalign{}
\endhead
\bottomrule\noalign{}
\endlastfoot
Add another scheduler & Implement \InlineCode{Job} and
\InlineCode{JobSubmitter}, then add a branch in \InlineCode{main}. \\
Publish URLs without a shared filesystem & Implement
\InlineCode{JobCommunicator} and its factory, and adapt launch scripts
to the new protocol. \\
Give a known model different HPC resources & Supply
\InlineCode{job\_<model\_name>.sh} in the selected script directory. \\
Change session lifetime & Extend the \InlineCode{UMBridge-Session}
request header and \InlineCode{/CloseSession} handling; preserve
one-operation behavior for clients without a session. \\
Bound startup time or detect failed jobs & Extend job-state inspection
and URL waiting with explicit failure/timeout handling. \\
Schedule based on configuration & Extend the manager/submitter interface
to receive the necessary configuration. \\
\end{longtable}

As written, the header has no include guard and contains
non-\InlineCode{inline} free function definitions. Its present
single-translation-unit use avoids some of the resulting problems; reuse
from multiple translation units would require addressing duplicate
inclusion/definition issues.

\hypertarget{reading-the-doxygen-reference}{%
\subsection{Reading the Doxygen
reference}\label{reading-the-doxygen-reference}}

The Doxygen comments alongside the declarations provide the API-level
reference: parameters, return values, ownership expectations, and
important failure conditions. This guide supplies the cross-class
control flow and the operational context needed to understand those
declarations together.

For an HTML reference, run Doxygen from \InlineCode{sandbox/C++/hpc}
with a configuration containing the following settings.
\InlineCode{EXTRACT\_PRIVATE} includes the environment conversion and
URL-filename helpers.

\begin{Shaded}
\begin{Highlighting}[]
\NormalTok{PROJECT\_NAME           = "UM{-}Bridge LoadBalancer"}
\NormalTok{OUTPUT\_DIRECTORY       = /tmp/loadbalancer{-}doxygen}
\NormalTok{INPUT                  = LoadBalancer.hpp LoadBalancer.cpp LoadBalancer.md}
\NormalTok{USE\_MDFILE\_AS\_MAINPAGE  = LoadBalancer.md}
\NormalTok{EXTRACT\_ALL            = YES}
\NormalTok{EXTRACT\_PRIVATE        = YES}
\NormalTok{GENERATE\_HTML          = YES}
\NormalTok{GENERATE\_LATEX         = NO}
\end{Highlighting}
\end{Shaded}

Save that configuration as a temporary Doxyfile, run
\InlineCode{doxygen\space{}<Doxyfile-path>}, and open
\InlineCode{/tmp/loadbalancer-doxygen/html/index.html}. Generating
documentation parses the sources; it does not require a successful
executable build or submit HPC jobs.

\end{document}
